\subsection*{A compact visual reader shows no reliable gain in Mandarin word order}
To isolate reading from drawing, we compare sentence images that contain
exactly the same glyphs in different orders. A new causal reader trained from
scratch receives continuous representations from a previously trained image
codec. Its eight-component writer and stop head score both complete orders.
This candidate contains 23.55 million parameters, including the frozen
7.42-million-parameter codec. During training it predicts future writing images
from Chinese book and question--answer sequences. The candidate and its
single-component control together complete 27,570 updates in 1,306 seconds on
one RTX 4090 D, reaching 1.416 GiB of peak allocated device memory. Earlier
codec fitting and corpus preparation contribute additional costs.

The reader at initialization prefers the grammatical order on 299 of 552
Mandarin pairs. After one epoch this count is 284 (51.4\%).
A visual bigram estimated on the same prediction targets scores 57.2\% with
half credit for ties. The five reader scores range from 48.0\% to 54.2\%
(Fig.~\ref{fig:visual-word-order}). Resampling whole constructions gives a
descriptive 95\% interval of 43.8--59.4\% for final accuracy and
$-19.6$ to $+14.7$ percentage points for its change from initialization.
These comparisons find no reliable gain in grammatical preference. They
assess the scoring of supplied images. Autonomous answer production requires
an additional decision about which image to draw at each step.
